% part: intuitionistic-logic
% chapter: introduction
% section: syntax

\documentclass[../../../include/open-logic-section]{subfiles}

\begin{document}

\olfileid{int}{int}{syn}

\olsection{স্বজ্ঞাবাদী যুক্তিবিদ্যার বাক্যবিন্যাস}

স্বজ্ঞাবাদী যুক্তিবিদ্যার বাক্যবিন্যাস বচনমূলক যুক্তিবিদ্যার মতোই।
ধ্রুপদি বচনমূলক যুক্তিবিদ্যায় একটি সংযোজককে অন্য সংযোজক দিয়ে
সংজ্ঞায়িত করা যায়; যেমন, $!A \lif !B$-কে $\lnot !A \lor !B$
দিয়ে, অথবা $!A \lor !B$-কে $\lnot(\lnot !A \land \lnot
!B)$ দিয়ে সংজ্ঞায়িত করা যায়। তাই ধ্রুপদি যুক্তিবিদ্যার বিবরণে
কোনো কোনো সংযোজককে এই সংজ্ঞাগুলির সংক্ষিপ্ত রূপ হিসেবে আনা হয়।
স্বজ্ঞাবাদী যুক্তিবিদ্যায় তা চলে না, দুটি ব্যতিক্রম ছাড়া:
$\lnot !A$-কে $!A \lif \lfalse$-এর সংক্ষিপ্ত রূপ হিসেবে
সংজ্ঞায়িত করা যায়---এবং প্রায়ই করা হয়। সেক্ষেত্রে অবশ্য
$\lfalse$-কে নিজে আর সংজ্ঞায়িত প্রতীক ধরা চলবে না!{} এছাড়া,
ধ্রুপদি যুক্তিবিদ্যার মতোই $!A \liff !B$-কে $(!A \lif !B) \land
(!B \lif !A)$ হিসেবে সংজ্ঞায়িত করা যায়।

বচনমূলক স্বজ্ঞাবাদী যুক্তিবিদ্যার !!^{formula}s গড়ে ওঠে
\emph{!!{propositional variable}s} এবং বচনমূলক ধ্রুবক~$\lfalse$
থেকে, \emph{যৌক্তিক সংযোজক} ব্যবহার করে। আমাদের আছে:

\begin{enumerate}
\item !!^a{denumerable} সেট~$\PVar$, যার !!{propositional variable}s
  হলো $\Obj p_0$, $\Obj p_1$, \dots
\item !!{falsity} নির্দেশকারী বচনমূলক ধ্রুবক~$\lfalse$।
\item যৌক্তিক সংযোজকগুলি: $\land$ (সংযোজন), $\lor$ (বিকল্প),
  $\lif$ (শর্তসূচক)।
\item বিরামচিহ্ন: (, ), এবং কমা।
\end{enumerate}

\begin{defn}[সূত্র]
\ollabel{defn:formulas} বচনমূলক স্বজ্ঞাবাদী যুক্তিবিদ্যার
\emph{!!{formula}s}-এর সেট~$\Frm[L_0]$ নিচের মতো আরোহক্রমে
সংজ্ঞায়িত:
\begin{enumerate}
\item $\lfalse$ একটি পারমাণবিক !!{formula}।

\item প্রত্যেক !!{propositional variable}~$\Obj p_i$ একটি পারমাণবিক
  !!{formula}।

\item $!A$ ও $!B$ !!{formula}s হলে $(!A \land
  !B)$ একটি !!a{formula}।

\item $!A$ ও $!B$ !!{formula}s হলে $(!A \lor !B)$
  একটি !!a{formula}।

\item $!A$ ও $!B$ !!{formula}s হলে $(!A \lif !B)$
  একটি !!a{formula}।

\tagitem{limitClause}{এগুলি ছাড়া আর কিছুই !!a{formula} নয়।}{}
\end{enumerate}
\end{defn}

উপরে আনা মৌলিক সংযোজকগুলির পাশাপাশি আমরা নিচের \emph{সংজ্ঞায়িত}
প্রতীকগুলিও ব্যবহার করি: $\lnot$ (নঞর্থকতা) এবং $\liff$
(!!{biconditional})। সংজ্ঞায়িত অপারেটর দিয়ে গঠিত সূত্রগুলি নিচের
অর্থে বুঝতে হবে:
\begin{enumerate}
\item $\lnot !A$ হলো $!A \lif \lfalse$-এর সংক্ষিপ্ত রূপ।

\item $!A \liff !B$ হলো $(!A \lif !B) \land (!B
  \lif !A)$-এর সংক্ষিপ্ত রূপ।
\end{enumerate}

$\lnot$-কে আনুষ্ঠানিকভাবে সংক্ষিপ্ত রূপ হিসেবে ধরা হলেও কখনো
কখনো সংজ্ঞায় সেটির জন্য আলাদা নিয়ম ও শর্ত দেব, যেন $\lnot$ মৌলিক
প্রতীক। প্রধানত অনুশীলনী দেওয়ার সুবিধার জন্যই তা করা হবে।

\end{document}
